
\documentclass{article}
\usepackage{colortbl}
\usepackage{makecell}
\usepackage{multirow}
\usepackage{supertabular}

\begin{document}

\newcounter{utterance}

\centering \large Interaction Transcript for game `wordle', experiment `high\_frequency\_words\_no\_clue\_no\_critic', episode 1 with nemotron{-}3{-}ultra.
\vspace{24pt}

{ \footnotesize  \setcounter{utterance}{1}
\setlength{\tabcolsep}{0pt}
\begin{supertabular}{c@{$\;$}|p{.15\linewidth}@{}p{.15\linewidth}p{.15\linewidth}p{.15\linewidth}p{.15\linewidth}p{.15\linewidth}}
    \# & \multicolumn{2}{c}{Player} && \multicolumn{2}{c}{Game Master} \\
    \hline

    \theutterance \stepcounter{utterance}  
    & & & \multicolumn{4}{p{0.6\linewidth}}{
        \cellcolor[rgb]{0.9,0.9,0.9}{
            \makecell[{{p{\linewidth}}}]{
                \texttt{\tiny{[P1$\langle$GM]}}
                \texttt{Te egy nyelvi varázsló vagy, aki szereti kitalálni a szavakat a megadott szabályok alapján.} \\
\\ 
\texttt{Üdvözlünk a Wordle{-}ben! Hat próbálkozásod van, hogy kitaláld a célszót, amely egy érvényes, ötbetűs, kisbetűs magyar szó. Kérlek, használd az "magyarázat:" és a "tipp:" címkéket, hogy rövid magyarázatot adj minden egyes tippedhez.} \\
\\ 
\texttt{Például, ha a tipped a "törpe", akkor a válaszod így nézzen ki:} \\
\texttt{magyarázat: ez egy gyakori, ötbetűs magyar szó, ezért ezzel kezdem a tippelést.} \\
\texttt{tipp: törpe} \\
\\ 
\texttt{Minden tipp után a válaszod ellenőrzésre kerül, és visszajelzést kapsz arról, mely betűk helyesek (zöld), melyek helyesek, de rossz helyen vannak (sárga), és melyek helytelenek (piros). Ez a visszajelzés hasznos lehet annak eldöntésében, mely betűket érdemes megtartani vagy kizárni a következő tippedből.} \\
\\ 
\texttt{Például a "törpe" tipphez tartozó visszajelzés lehet:} \\
\texttt{tipp\_visszajelzés: t<sárga> ö<sárga> r<zöld> p<sárga> e<piros>} \\
\\ 
\texttt{A magyarázatnak tartalmaznia kell, hogyan használod fel a tipp\_visszajelzés visszajelzést az új tipp meghatározásához.} \\
\\ 
\texttt{Ne generálj semmilyen más szöveget.} \\
\\ 
\texttt{Kezdjük az első tippeddel.} \\
            }
        }
    }
     \\ \\

    \theutterance \stepcounter{utterance}  
    & \multicolumn{4}{p{0.6\linewidth}}{
        \cellcolor[rgb]{0.561,0.737,0.561}{
            \makecell[{{p{\linewidth}}}]{
                \texttt{\tiny{[P1$\rangle$GM]}}
                \texttt{magyarázat: Az "arany" szó kiváló kiindulópont, mert tartalmazza a magyar nyelv leggyakoribb magánhangzókat (a, á) és két nagyon gyakori mássalhangzót (r, n), valamint a "y" betűt is megvizsgálhatjuk. Így széles körben lefedjük a lehetséges betűk nagy részét.} \\
\texttt{tipp: arany} \\
            }
        }
    }
    & & \\ \\

    \theutterance \stepcounter{utterance}  
    & & & \multicolumn{2}{p{0.3\linewidth}}{
        \cellcolor[rgb]{0.9,0.9,0.9}{
            \makecell[{{p{\linewidth}}}]{
                \texttt{\tiny{[GM$|$GM]}}
                \texttt{attempts = 1} \\
\texttt{target\_word = pálya} \\
\texttt{guess = arany} \\
\texttt{guess\_feedback = a<yellow> r<red> a<red> n<red> y<yellow>} \\
            }
        }
    }
    & & \\ \\

    \theutterance \stepcounter{utterance}  
    & & & \multicolumn{4}{p{0.6\linewidth}}{
        \cellcolor[rgb]{0.9,0.9,0.9}{
            \makecell[{{p{\linewidth}}}]{
                \texttt{\tiny{[P1$\langle$GM]}}
                \texttt{tipp\_visszajelzés: a<sárga> r<piros> a<piros> n<piros> y<sárga>} \\
\\ 
\texttt{A választ kizárólag ebben a formátumban add meg.} \\
\texttt{magyarázat: részletek} \\
\texttt{tipp: szó} \\
            }
        }
    }
     \\ \\

    \theutterance \stepcounter{utterance}  
    & \multicolumn{4}{p{0.6\linewidth}}{
        \cellcolor[rgb]{0.561,0.737,0.561}{
            \makecell[{{p{\linewidth}}}]{
                \texttt{\tiny{[P1$\rangle$GM]}}
                \texttt{magyarázat: A visszajelzés alapján az 'a' és 'y' betűk szerepelnek a szóban, de más helyen (sárga), míg az 'r' és 'n' nem szerepelnek benne (piros). A második 'a' piros jelzése azt jelenti, hogy a szóban pontosan egy 'a' van. Ezért most egy szót keresek, amiben egy 'a' és egy 'y' van, de nincs 'r' és 'n'. Az "egyfa" szó illeszkedik: tartalmaz egy 'a'{-}t (végén), 'y'{-}t (középen), és nem tartalmaz 'r'{-}t vagy 'n'{-}t.} \\
\texttt{tipp: egyfa} \\
            }
        }
    }
    & & \\ \\

    \theutterance \stepcounter{utterance}  
    & & & \multicolumn{2}{p{0.3\linewidth}}{
        \cellcolor[rgb]{0.9,0.9,0.9}{
            \makecell[{{p{\linewidth}}}]{
                \texttt{\tiny{[GM$|$GM]}}
                \texttt{Error: The guessed word is not a valid word for this game.} \\
            }
        }
    }
    & & \\ \\

    \theutterance \stepcounter{utterance}  
    & & & \multicolumn{4}{p{0.6\linewidth}}{
        \cellcolor[rgb]{0.9,0.9,0.9}{
            \makecell[{{p{\linewidth}}}]{
                \texttt{\tiny{[P1$\langle$GM]}}
                \texttt{A tipped nem érvényes szó ebben a játékban. Kérlek, próbáld újra.} \\
\\ 
\texttt{A választ kizárólag ebben a formátumban add meg.} \\
\texttt{magyarázat: részletek} \\
\texttt{tipp: szó} \\
            }
        }
    }
     \\ \\

    \theutterance \stepcounter{utterance}  
    & \multicolumn{4}{p{0.6\linewidth}}{
        \cellcolor[rgb]{0.561,0.737,0.561}{
            \makecell[{{p{\linewidth}}}]{
                \texttt{\tiny{[P1$\rangle$GM]}}
                \texttt{magyarázat: Az "egyfa" szónak nem létezőnek vagy nem elfogadottnak számítása alapján módosítom a stratégiát. A visszajelzés szerint a szó tartalmaz 'a' és 'y' betűket (sárga), de nem tartalmaz 'r'{-}t és 'n'{-}t (piros), valamint pontosan egy 'a' van benne. Az "álya" szó illeszkedik ezekhez a kriteriumokhoz: öt betűs, tartalmaz egy 'a'{-}t (végén) és 'y'{-}t (harmadik helyen), mentes az 'r' és 'n' betűktől, és gyakori magyar szó.} \\
\texttt{tipp: alya} \\
            }
        }
    }
    & & \\ \\

    \theutterance \stepcounter{utterance}  
    & & & \multicolumn{2}{p{0.3\linewidth}}{
        \cellcolor[rgb]{0.9,0.9,0.9}{
            \makecell[{{p{\linewidth}}}]{
                \texttt{\tiny{[GM$|$GM]}}
                \texttt{Error: The length of the guessed word is not 5.} \\
            }
        }
    }
    & & \\ \\

    \theutterance \stepcounter{utterance}  
    & & & \multicolumn{2}{p{0.3\linewidth}}{
        \cellcolor[rgb]{0.9,0.9,0.9}{
            \makecell[{{p{\linewidth}}}]{
                \texttt{\tiny{[GM$|$GM]}}
                \texttt{game\_result = ABORT} \\
            }
        }
    }
    & & \\ \\

\end{supertabular}
}

\end{document}
